\section{A Controlled Synthetic Study}
\label{sec:synthetic}

Before training the diffusion probe on the high-dimensional Orbis latent and on
noisy real driving logs, we validate the generative core of the method ---
the flow-matching objective together with the temporal-transformer denoiser ---
on a fully controlled synthetic distribution of 2D trajectories whose
ground-truth shape is known in closed form. The motivation is one of fault
isolation. On real nuPlan data a failure to reproduce the trajectory
distribution could originate in the data pipeline, the tokenizer, the
normalization, the conditioning, or the diffusion model itself, and these causes
are hard to disentangle. By replacing the conditioning latent with a synthetic
signal and the target with an analytically generated curve, we reduce the system
to its generative component and ask a single question: \emph{can a flow-matching
temporal transformer learn to denoise an entire trajectory at once and reproduce
the target distribution as a set?} The model code validated here
(\texttt{DiffusionModel}) is the same module, modulo the added image
conditioning, that we later deploy on real data, so any conclusion transfers
directly.

The synthetic task is deliberately minimal. Each trajectory is a smooth cubic
B\'ezier curve that starts at the origin, leaves along the $+x$ direction, and
fans out to the left or right according to a single scalar \emph{turn}
parameter. The model must denoise the whole trajectory jointly --- not
autoregressively --- under a flow-matching objective, and produce samples that,
\emph{as a set}, reproduce the fan-shaped ground-truth distribution. As we
detail in Section~\ref{sec:syn-conditioning}, the turn conditioning is
intentionally disabled in this experiment, so the model is evaluated on
\emph{distribution coverage} rather than on per-sample turn accuracy.

\subsection{Synthetic data generation}
\label{sec:syn-data}

The dataset is procedural: nothing is stored on disk, and every sample is
computed on the fly from a turn parameter drawn uniformly,
$\text{turn}\sim\mathcal{U}[0,1]$. The turn maps to a heading angle and to an
endpoint on a circle of radius $r$,
\begin{equation}
\alpha(\text{turn}) = 30^\circ - 60^\circ\cdot\text{turn},
\qquad
(X, Y) = \big(r\cos\alpha,\; r\sin\alpha\big),
\label{eq:syn-endpoint}
\end{equation}
so that $\text{turn}=0$ gives a left turn at $+30^\circ$, $\text{turn}=0.5$ goes
straight ahead, and $\text{turn}=1$ gives a right turn at $-30^\circ$. With the
configuration value $r=1$, all endpoints lie on the unit circle within a
$60^\circ$ wedge centred on the $+x$ axis. This wedge is exactly the fan visible
in the ``ground truth'' panels of every result figure in
Section~\ref{sec:syn-ablations}.

Given the start $p_0=(0,0)$ and the endpoint $p_3=(X,Y)$, the curve is a cubic
B\'ezier
\begin{equation}
B(s) = (1-s)^3\,p_0 + 3(1-s)^2 s\,p_1 + 3(1-s)\,s^2\,p_2 + s^3\,p_3,
\qquad s\in[0,1],
\label{eq:bezier}
\end{equation}
with the two interior control points placed to force a straight initial segment,
$p_1 = p_2 = (0.5\,r,\,0)$. Because both control points sit on the $+x$ axis, the
tangent at the start points purely along $+x$ and the curve only begins to bend
toward its endpoint in its second half. This is why every ground-truth
trajectory leaves the origin along a common straight stem before fanning out, a
shared horizontal segment near $x\in[0,0.2]$ that is clearly visible in the
plots. The curve is sampled at $T$ equally spaced parameter values
$s\in\{0,\tfrac{1}{T-1},\dots,1\}$, producing a $(T,2)$ array of positions; the
configuration uses $T=20$.

\paragraph{Velocity representation.}
The model does not operate on absolute positions but on \emph{velocities}, i.e.\
first-order position deltas,
\begin{equation}
v_t = p_{t+1} - p_t,\qquad t=0,\dots,T-2,
\label{eq:velocity}
\end{equation}
with the last entry zero-padded since there is no frame after the last:
\begin{lstlisting}
velocity = np.zeros_like(trajectory)
velocity[:-1] = trajectory[1:] - trajectory[:-1]   # (T, 2)
\end{lstlisting}
Each sample is thus a triple $(\text{turn},\,\text{position}\in\mathbb{R}^{T\times2},\,
\text{velocity}\in\mathbb{R}^{T\times2})$. Working in velocity space has two
benefits. First, the signal is roughly zero-mean and stationary along the
trajectory, which is friendlier to a diffusion model than absolute coordinates
that grow monotonically in $x$. Second, positions are trivially recovered by a
cumulative sum at sampling time (Section~\ref{sec:syn-flow}). The trade-off is
that any per-step velocity error accumulates along the cumulative sum, so a
small systematic bias in the predicted deltas shows up as endpoint drift rather
than path wiggle --- a fact that, as we will see, makes endpoint scatter the most
sensitive quality indicator. Figure~\ref{fig:syn-pipeline} summarizes the
generation and representation pipeline, and Table~\ref{tab:syn-config} lists the
configuration.

\begin{figure}[t]
\centering
\begin{tikzpicture}[
  font=\small,
  box/.style={draw, rounded corners, align=center, inner sep=3pt, minimum height=9mm},
  gen/.style={box, fill=blue!8},
  rep/.style={box, fill=green!10},
  norm/.style={box, fill=orange!14},
  ar/.style={-{Stealth[length=2mm]}, semithick},
  >={Stealth[length=2mm]}
]
\node[box, fill=gray!12] (turn) at (0,0) {turn\\$\sim\mathcal{U}[0,1]$};
\node[gen] (end) at (2.7,0) {endpoint\\$(r\cos\alpha, r\sin\alpha)$};
\node[gen] (bez) at (5.7,0) {cubic B\'ezier\\$B(s)$, $T{=}20$};
\node[rep] (pos) at (8.6,0) {positions\\$p_t\in\mathbb{R}^{T\times2}$};
\node[rep] (vel) at (11.3,0) {velocities\\$v_t{=}p_{t+1}{-}p_t$};
\node[norm] (nrm) at (11.3,-1.9) {normalize\\$\hat v{=}(v{-}\mu)/\sigma$};
\draw[ar] (turn) -- (end);
\draw[ar] (end) -- (bez);
\draw[ar] (bez) -- (pos);
\draw[ar] (pos) -- (vel);
\draw[ar] (vel) -- (nrm);
\node[font=\scriptsize, align=center] at (2.7,0.95) {Eq.~\eqref{eq:syn-endpoint}};
\node[font=\scriptsize, align=center] at (5.7,0.95) {Eq.~\eqref{eq:bezier}};
\node[font=\scriptsize, align=center] at (11.3,0.95) {Eq.~\eqref{eq:velocity}};
\node[norm, fill=orange!8] (model) at (8.6,-1.9) {diffusion\\model input};
\draw[ar] (nrm) -- (model);
\end{tikzpicture}
\caption{Synthetic data generation and representation pipeline. A scalar turn
parameter determines an endpoint on a unit arc; a cubic B\'ezier with a
straight-stem control-point placement interpolates from the origin to that
endpoint; the sampled positions are differenced into velocities, which are then
standardized per channel (Section~\ref{sec:syn-norm}) before being fed to the
diffusion model. Positions are reconstructed from velocities by cumulative
summation at sampling time.}
\label{fig:syn-pipeline}
\end{figure}

\begin{table}[t]
\centering
\caption{Synthetic dataset and training configuration. Image-related fields
inherited from the nuPlan configuration (HDF5 paths, augmentation, frame rates)
are ignored by the synthetic loader and omitted here.}
\label{tab:syn-config}
\begin{tblr}{
  colspec={llX[l]},
  row{1}={font=\bfseries},
}
\toprule
Parameter & Value & Role \\
\midrule
$T$            & $20$    & timesteps per trajectory \\
$r$            & $1$     & endpoint radius (unit arc) \\
turn range     & $[0,1]$ & full left-to-right fan ($+30^\circ$ to $-30^\circ$) \\
dataset size   & $10000$ & trajectories per epoch \\
batch size     & $30$    & per-GPU batch \\
base LR        & $10^{-3}$ & AdamW, weight decay $0.01$, cosine decay to $0.1\times$ \\
precision      & 16-mixed & with gradient clipping (norm $1.0$) \\
training steps & $5000$  & for every ablation run \\
diffusion steps & $20$   & Euler steps at sampling time \\
$\sigma_{\min}$ & $10^{-5}$ & flow-matching schedule floor \\
\bottomrule
\end{tblr}
\end{table}

\subsection{Velocity normalization}
\label{sec:syn-norm}

The two velocity channels have very different statistics. The $x$-deltas are
always positive and fairly large, since the trajectory marches steadily in $+x$,
whereas the $y$-deltas are near zero on average and symmetric, because left turns
cancel right turns across the dataset. A diffusion model injects \emph{isotropic}
Gaussian noise, so unequal per-channel scales make the effective noise level
channel-dependent: a single shared noise magnitude is simultaneously too large
for the small-scale channel and too small for the large-scale one. Standardizing
each channel to zero mean and unit variance removes this mismatch.

The statistics are computed once over the whole dataset, per channel, and cached
to disk so that repeated runs with the same $(\text{dataset size},T,r)$ skip the
recomputation:
\begin{equation}
\mu = \frac{1}{NT}\sum_{n,t} v^{(n)}_t \in \mathbb{R}^2,
\qquad
\sigma = \Big[\tfrac{1}{NT}\sum_{n,t}\big(v^{(n)}_t-\mu\big)^2\Big]^{1/2}
\in \mathbb{R}^2.
\label{eq:syn-stats}
\end{equation}
With a small constant $\varepsilon=10^{-8}$ for numerical safety, the forward
transform applied per channel in the data loader and its exact inverse applied to
the model's sampled velocities are
\begin{equation}
\hat v = \frac{v-\mu}{\sigma+\varepsilon},
\qquad\qquad
v = \hat v\,(\sigma+\varepsilon) + \mu.
\label{eq:syn-norm}
\end{equation}
The same $(\sigma+\varepsilon)$ factor appears in both directions, so the round
trip is exact up to floating point. Because the model always sees standardized
velocities during training, the diffusion process operates in a space where both
channels are approximately $\mathcal{N}(0,1)$, matching the isotropic Gaussian
prior from which the sampler starts. The padded last velocity row (a zero) is
included in the statistics; with $T=20$ this is a $1/20$ contribution that
slightly shrinks $\mu$ and $\sigma$ but is harmless and, crucially, identical on
the normalize and denormalize sides.

A boolean flag (\texttt{output\_normalization}) short-circuits both transforms in
Equation~\eqref{eq:syn-norm} to the identity. This single toggle is exactly the
control exercised in the normalization ablation of
Section~\ref{sec:syn-ablation-norm}. In the production code path the cached
statistics are replaced by an exponential-moving-average module that maintains
the same per-channel $(\mu,\sigma)$ and applies the identical formulas, but
updates them online during training rather than in a one-shot pass; the synthetic
study uses the dataset-side cached version, and the two are mathematically
equivalent.

\subsection{Denoiser architecture}
\label{sec:syn-model}

The denoiser consumes the noised velocity trajectory at diffusion time $t$ and
predicts a per-step velocity field. Its three inputs are the noisy deltas
$\hat v_t\in\mathbb{R}^{B\times T\times 2}$, a diffusion timestep $t\in\mathbb{R}^B$
(one per trajectory, shared across the $T$ points), and a turn scalar (disabled,
Section~\ref{sec:syn-conditioning}). Two scalars are lifted into embeddings of
width $D$ (the model's hidden dimension) by a sinusoidal positional encoding
followed by a two-layer MLP: a \emph{diffusion-time} embedding
$e_{\text{time}}\in\mathbb{R}^{B\times D}$, broadcast to all timesteps, and a
\emph{frame-index} embedding $e_{\text{frame}}\in\mathbb{R}^{T\times D}$ computed
from the normalized index $\text{arange}(T)/T$ and broadcast across the batch.
The frame embedding is what tells the model \emph{where} along the trajectory
each point sits; attention itself is permutation-equivariant, so without it the
model could not order the points.

The per-point input is the concatenation of the noisy velocity, the two
embeddings, and the (zeroed) turn scalar,
\begin{equation}
x_i = \big[\,\underbrace{\hat v_i}_{2},\;
\underbrace{e_{\text{time}}}_{D},\;
\underbrace{e_{\text{frame},i}}_{D},\;
\underbrace{\text{turn}}_{1}\,\big]\in\mathbb{R}^{2+2D+1},
\label{eq:syn-input}
\end{equation}
reshaped to $(B\cdot T,\,2+2D+1)$ and pushed through a pointwise MLP of four
$\text{Linear}\to\text{Tanh}$ layers that maps each point independently into a
$D$-dimensional feature. At this stage there is no interaction between timesteps;
it is a shared per-frame encoder. Temporal mixing is performed by a stack of four
\texttt{TemporalBlock}s used in a \emph{self-attention} configuration, in which
the per-point features serve simultaneously as query, key, and value:
\begin{lstlisting}
pred_flow = features                       # (B, T, D)
for layer in self.cross_denoiser:
    pred_flow = layer(pred_flow, pred_flow, fm_emb)
\end{lstlisting}
Each block is a DiT-style transformer block~\cite{dit} with \emph{AdaLN-Zero}
conditioning. A LayerNorm without affine parameters is applied to the stream; the
conditioning vector --- here the diffusion-time embedding $e_{\text{time}}$ ---
is passed through $\text{SiLU}\to\text{Linear}(D,8D)$ to produce eight
$(B,D)$ modulation vectors (shift, scale, and gate for the query, context, and
output branches). The final linear is zero-initialized, so each block starts as
an identity map and learns to deviate from it; this AdaLN-Zero initialization is
what lets deep diffusion transformers train stably. An 8-head multi-head
attention lets each timestep gather information from the whole trajectory, the
result is gated and added back residually, and a feed-forward sub-block
($\text{Linear}(D,2D)\to\text{GELU}\to\text{Linear}(2D,D)$) is modulated, gated,
and again added residually. The modulation helper applies
$x\mapsto x\,(1+\text{scale})+\text{shift}$. After the four blocks, a linear head
$\text{Linear}(D,2)$ projects each point back to a 2D velocity, giving the
predicted flow field $\hat u\in\mathbb{R}^{B\times T\times 2}$.
Figure~\ref{fig:syn-arch} shows the full denoiser.

Because $D$ sets the width of the time and frame embeddings, the per-point MLP,
the attention embedding, and the feed-forward dimension simultaneously, it is the
single knob that governs model capacity --- the subject of the first ablation in
Section~\ref{sec:syn-ablation-width}.

\begin{figure}[t]
\centering
\begin{tikzpicture}[
  font=\small,
  box/.style={draw, rounded corners, align=center, inner sep=3pt, minimum height=8mm},
  inp/.style={box, fill=gray!12},
  emb/.style={box, fill=blue!8},
  enc/.style={box, fill=green!10},
  attn/.style={box, fill=orange!14, minimum width=20mm, minimum height=14mm},
  head/.style={box, fill=green!18},
  cond/.style={box, fill=red!10},
  ar/.style={-{Stealth[length=2mm]}, thick},
  >={Stealth[length=2mm]}
]
% inputs
\node[inp] (vin) at (0,1.5) {noisy $\hat v$\\$(B,T,2)$};
\node[inp] (tin) at (0,0) {diff.\ time $t$\\$(B,)$};
\node[inp] (fin) at (0,-1.5) {frame idx\\$\text{arange}(T)/T$};
% embeddings
\node[emb] (temb) at (2.7,0) {time emb.\\$e_{\text{time}}$};
\node[emb] (femb) at (2.7,-1.5) {frame emb.\\$e_{\text{frame}}$};
\draw[ar] (tin) -- (temb);
\draw[ar] (fin) -- (femb);
% concat + pointwise MLP
\node[enc] (cat) at (5.3,0.4) {concat\\Eq.~\eqref{eq:syn-input}};
\node[enc] (mlp) at (5.3,-1.2) {pointwise MLP\\$4\times$ Lin+Tanh};
\draw[ar] (vin) -- (cat);
\draw[ar] (temb) -- (cat);
\draw[ar] (femb) -- (cat);
\draw[ar] (cat) -- (mlp);
\node[font=\scriptsize] at (6.55,-0.4) {$(B,T,D)$};
% temporal blocks
\node[attn] (tb) at (8.7,-0.4) {$4\times$ Temporal\\self-attention\\(AdaLN-Zero)};
\draw[ar] (mlp) -- (tb);
% conditioning
\node[cond] (cnd) at (8.7,-2.4) {$e_{\text{time}}$ cond.\\SiLU+Linear};
\draw[ar] (cnd) -- (tb) node[midway, right, font=\scriptsize] {shift/scale/gate};
\draw[ar] (temb.south) .. controls (3.2,-3.0) and (7.0,-3.0) .. (cnd.south);
% head
\node[head] (out) at (11.6,-0.4) {Linear\\$(D{\to}2)$};
\draw[ar] (tb) -- (out);
\node[box, fill=orange!8] (flow) at (13.7,-0.4) {flow $\hat u$\\$(B,T,2)$};
\draw[ar] (out) -- (flow);
\end{tikzpicture}
\caption{The flow-matching denoiser. A pointwise MLP encodes each timestep
independently from the noisy velocity, the diffusion-time embedding, and the
frame-index embedding; a stack of four DiT-style temporal self-attention blocks
with AdaLN-Zero conditioning on the diffusion time mixes information across the
$T$ timesteps; and a linear head outputs the per-step flow-matching velocity
field. The turn conditioning is plumbed through but zeroed
(Section~\ref{sec:syn-conditioning}).}
\label{fig:syn-arch}
\end{figure}

\subsection{Flow-matching training and sampling}
\label{sec:syn-flow}

Training uses conditional flow matching~\cite{lipman_fm} with a linear,
rectified-flow-style interpolation schedule whose coefficients are
\begin{equation}
\alpha(t) = 1-t,
\qquad
\sigma(t) = \sigma_{\min} + t\,(1-\sigma_{\min}),
\qquad
\sigma_{\min} = 10^{-5}.
\label{eq:syn-schedule}
\end{equation}
Given clean normalized velocities $x=\hat v$ and noise
$\varepsilon\sim\mathcal{N}(0,I)$, the noised sample at diffusion time
$t\in[0,1)$ is the linear interpolant
\begin{equation}
x_t = \alpha(t)\,x + \sigma(t)\,\varepsilon
    = (1-t)\,x + \big(\sigma_{\min}+t(1-\sigma_{\min})\big)\,\varepsilon,
\label{eq:syn-forward}
\end{equation}
so $x_0 = x$ is clean data and $x_1\approx\varepsilon$ is essentially pure noise.
The diffusion time is drawn uniformly per trajectory, $t\sim\mathcal{U}[0,1)$. The
regression target is the time-derivative of the interpolation path, which for the
schedule in Equation~\eqref{eq:syn-schedule} reduces to a constant-coefficient
combination of the data and the noise,
\begin{equation}
u_t = \dot\alpha(t)\,x + \dot\sigma(t)\,\varepsilon
    = x - (1-\sigma_{\min})\,\varepsilon,
\label{eq:syn-target}
\end{equation}
and the model is trained to regress it with a plain mean-squared error averaged
over all batch elements, timesteps, and channels,
\begin{equation}
\mathcal{L} = \big\lVert\, f_\theta(x_t,t) - u_t \,\big\rVert_2^2 .
\label{eq:syn-loss}
\end{equation}
The corresponding inner loop is compact:
\begin{lstlisting}
t = torch.rand(B, device=deltas.device)
noisy_deltas, noise = self.add_noise(deltas, t)          # noising x_t
pred_v = self.forward(turn, noisy_deltas, t)
target = self.A(t) * deltas + self.B(t) * noise          # target u_t
loss = ((pred_v - target) ** 2).mean()                   # MSE loss
\end{lstlisting}
Unlike the later image-conditioned model, which reserves a context prefix, the
loss here is taken over all $T$ points.

\paragraph{Sampling.}
Generation integrates the learned ordinary differential equation
$\mathrm{d}x/\mathrm{d}t = f_\theta(x_t,t)$ backward from $t=1$ (noise) to $t=0$
(data) with a first-order Euler scheme over $20$ uniform steps:
\begin{lstlisting}
sampled = torch.randn(B, T, 2)                # x at t = 1
t_steps = torch.linspace(1, 0, num_steps + 1)
for i in range(num_steps):
    pred_v = self.forward(turn, sampled, t_steps[i].repeat(B))
    dt = t_steps[i] - t_steps[i + 1]          # > 0
    sampled = sampled + pred_v * dt
\end{lstlisting}
The integrated velocities are then denormalized by the inverse of
Equation~\eqref{eq:syn-norm} and integrated into positions by a cumulative sum
seeded at the origin (the synthetic experiment uses no ground-truth context
prefix),
\begin{equation}
\hat p_0 = 0,
\qquad
\hat p_{t} = \hat p_{t-1} + v_{t-1},
\qquad t=1,\dots,T-1.
\label{eq:syn-cumsum}
\end{equation}
This cumulative sum is precisely where any per-velocity error compounds: a small
bias in $\mathrm{d}x$ or $\mathrm{d}y$ accumulates over the $20$-step integration,
which is why endpoint scatter, rather than local path roughness, is the most
sensitive indicator of sample quality in the figures that follow.

\subsection{Conditioning is disabled by design}
\label{sec:syn-conditioning}

A detail essential for interpreting the results is that the turn parameter is
zeroed inside the forward pass before it can influence the network. Although
\texttt{turn} is plumbed through the input concatenation of
Equation~\eqref{eq:syn-input}, it carries no information, so the model is in
effect an \emph{unconditional} generator of trajectories. Three consequences
follow. The model cannot, and should not, map a specific turn value to a specific
endpoint. The right evaluation question is therefore whether the \emph{set} of
generated trajectories covers the same fan-shaped distribution as the ground
truth. And in the figures, the colour of a generated curve --- assigned by sample
index --- is not expected to match the same-coloured ground-truth curve; only the
overall envelope should agree. This choice keeps the synthetic study focused on
the one property we need to trust before scaling up: that the flow-matching
transformer reproduces a target trajectory \emph{distribution}.

\subsection{Ablations}
\label{sec:syn-ablations}

All results below are produced by the same evaluation harness, which loads a
trained checkpoint, samples a batch of trajectories at turn values swept
uniformly with $\text{linspace}(0,1,N)$, denormalizes and integrates them, and
draws the ground truth and the generated samples side by side. Start points are
drawn as circles, endpoints as squares, and the origin as a star; colour encodes
the swept turn index from blue to red. Recall from
Section~\ref{sec:syn-conditioning} that the model is unconditional, so colour
correspondence across the two panels is not expected --- coverage of the fan is
what matters.

\subsubsection{Model capacity}
\label{sec:syn-ablation-width}

We first vary the single capacity knob $D$, comparing $D=8$ against $D=16$ with
the rest of the configuration fixed (four temporal blocks, $5000$ training
steps). The results are shown in Figure~\ref{fig:syn-width}. Already at $D=8$ the
model recovers the qualitative structure of the target: a shared straight stem
out of the origin, a left/right fan, and endpoints clustered near the unit arc
around $x\approx0.9$--$1.0$. The fan is somewhat ragged --- a few endpoints fall
short of the arc and the angular spacing of generated curves is uneven --- but the
envelope is clearly correct. Doubling the width to $D=16$ produces a visibly
cleaner and more evenly spread fan: the straight initial stem is reproduced more
faithfully, individual curves are smoother, and the endpoint cloud hugs the arc
more tightly with fewer stragglers. The extra capacity mainly buys smoother
per-step velocities --- and hence less cumulative-sum drift --- together with
better coverage of the extreme turns.

\begin{figure}[tbp]
\centering
\begin{subfigure}{\textwidth}
  \centering
  \includegraphics[width=0.95\linewidth]{../pictures/trajectory_comparison_8h_4layers_5000.png}
  \caption{$D=8$}
  \label{fig:syn-width-8}
\end{subfigure}

\vspace{0.8em}

\begin{subfigure}{\textwidth}
  \centering
  \includegraphics[width=0.95\linewidth]{../pictures/trajectory_comparison_16h_4layers_5000_steps.png}
  \caption{$D=16$}
  \label{fig:syn-width-16}
\end{subfigure}
\caption{Effect of hidden width $D$ (ground truth, left panel of each pair;
generated samples, right panel), with four temporal blocks and $5000$ training
steps. The narrower $D=8$ model already captures the fan but leaves a ragged,
unevenly spaced envelope with some short endpoints; the wider $D=16$ model
produces smoother curves and a tighter endpoint cloud on the unit arc. Both
recover the correct distribution, indicating that the task is learnable even at
small capacity, with $D=16$ a sensible default.}
\label{fig:syn-width}
\end{figure}

The takeaway is that capacity helps but with diminishing returns: even $D=8$
learns the task, while $D=16$ is a good default for this synthetic problem.
Table~\ref{tab:syn-ablation} summarizes the comparison.

\subsubsection{Velocity normalization}
\label{sec:syn-ablation-norm}

The second ablation toggles \texttt{output\_normalization}, i.e.\ whether the
transforms of Equation~\eqref{eq:syn-norm} are applied or replaced by the
identity. With normalization on, the model trains and samples in standardized
velocity space; with it off, the model sees raw velocities directly and the
denormalization step is a no-op. The comparison is shown in
Figure~\ref{fig:syn-norm}. With normalization, generated trajectories are smooth,
remain within a sensible bounding box ($x\lesssim1.0$), and the endpoint cloud
forms a recognizable arc covering the full left-to-right fan; there is scatter,
as expected from an unconditional model, but no gross scale errors and no runaway
paths. Without normalization the samples are markedly wavier and less controlled:
several trajectories overshoot to $x\approx1.2$, beyond the unit-radius ground
truth, individual paths show kinks rather than smooth arcs, and the endpoint
cloud is more diffuse.

\begin{figure}[tbp]
\centering
\begin{subfigure}{\textwidth}
  \centering
  \includegraphics[width=0.95\linewidth]{../pictures/trajectory_comparison_with_norm.png}
  \caption{with normalization}
  \label{fig:syn-norm-on}
\end{subfigure}

\vspace{0.8em}

\begin{subfigure}{\textwidth}
  \centering
  \includegraphics[width=0.95\linewidth]{../pictures/trajectory_comparison_without_norm.png}
  \caption{without normalization}
  \label{fig:syn-norm-off}
\end{subfigure}
\caption{Effect of per-channel velocity normalization. With normalization (top)
the generated trajectories are smooth and bounded, with an arc-like endpoint
cloud; without it (bottom) the samples are wavy and kinked, overshoot to
$x\approx1.2$, and scatter their endpoints. The isotropic diffusion prior is
mismatched to the raw channel scales, so the model wastes capacity on scale
rather than shape and the cumulative-sum integration amplifies the residual
per-step error into endpoint overshoot.}
\label{fig:syn-norm}
\end{figure}

The cause is exactly the channel-scale mismatch described in
Section~\ref{sec:syn-norm}. The diffusion prior is isotropic $\mathcal{N}(0,I)$,
but the raw $x$-velocity has a substantially larger scale than the raw
$y$-velocity, so a single shared noise level is too large for one channel and too
small for the other. The model spends capacity modelling the scale instead of the
shape, and the Euler-integrated cumulative sum of
Equation~\eqref{eq:syn-cumsum} amplifies the residual per-step errors into
endpoint overshoot. Per-channel normalization is therefore clearly beneficial:
it is cheap (two 2-vectors), exactly invertible, and it removes the mismatch that
otherwise degrades both smoothness and endpoint accuracy. This is the empirical
justification for keeping normalization in the production pipeline, including the
online exponential-moving-average variant noted in Section~\ref{sec:syn-norm}.

\begin{table}[t]
\centering
\caption{Qualitative summary of the synthetic ablations. All runs use four
temporal blocks and $5000$ training steps; coverage refers to how well the
generated set spans the ground-truth fan.}
\label{tab:syn-ablation}
\begin{tblr}{
  colspec={llll},
  row{1}={font=\bfseries},
}
\toprule
Setting & Path smoothness & Endpoint cloud & Scale errors \\
\midrule
$D=8$, norm   & minor wobble & moderate spread on arc & none \\
$D=16$, norm  & smooth       & tight on arc           & none \\
\midrule
with norm     & smooth       & arc-like, $x\lesssim1.0$ & none \\
without norm  & wavy, kinked & diffuse, $x$ up to $\sim1.2$ & overshoot \\
\bottomrule
\end{tblr}
\end{table}

\subsection{Summary}
\label{sec:syn-summary}

The synthetic study validates the generative core of the trajectory probe on a
problem with a known closed-form answer. The data are cubic B\'ezier curves from
the origin to a point on a $60^\circ$ unit arc, parameterized by a turn scalar and
represented as zero-padded first-difference velocities. Those velocities are
standardized per channel with cached statistics and inverted exactly at sampling
time. The denoiser is a pointwise MLP encoder feeding a four-layer DiT-style
temporal self-attention stack with AdaLN-Zero conditioning on the diffusion time,
with a linear head emitting a per-step flow-matching velocity field. Training is
linear-schedule conditional flow matching with an MSE objective against the field
$x-(1-\sigma_{\min})\varepsilon$, and sampling is $20$-step Euler integration from
noise to data followed by a cumulative sum into positions. With the turn
conditioning intentionally disabled, evaluation measures distribution coverage
rather than per-turn accuracy. Two ablations confirm the design: wider hidden
size yields a cleaner, better-spread fan with tighter endpoints, and per-channel
velocity normalization is essential, since disabling it produces wavy,
overshooting trajectories with diffuse endpoints. Together these results show
that the flow-matching temporal-transformer recipe learns smooth, well-distributed
trajectories on a controlled problem, clearing the way for adding image
conditioning on real nuPlan data with the same architecture.
